\clearpage
\section{Selecting primary and secondary chains}
\label{sec:chaining-primary-secondary}

Several chains may match the same part of a query at different target positions.
Minimap2 considers them from highest to lowest chaining score. A chain becomes
\emph{secondary} when an already accepted \emph{primary} covers enough of the
same query interval. Otherwise it becomes another primary. This allows
different parts of a query to map to different places, as in a split
alignment~\cite{minimap2_paper,minimap2_software}.

Let \(I_c=[q_c^{\mathrm{start}},q_c^{\mathrm{end}})\) be the query interval
of chain \(c\), and \(s_c\) its score. In
Algorithm~\ref{alg:chaining-primary-secondary}, \(|I|\) counts query bases.
The value \(u\) counts bases not covered by \emph{any} accepted primary.
Subtracting this contribution makes a chain with substantial new coverage
less likely to become secondary. The default overlap threshold is
\(\lambda=0.5\). The uncovered-length limit is \(L=2^{31}-1\), so it
normally adds no restriction.

\begin{algorithm}[H]
\caption{Minimap2 primary and secondary chain selection}
\label{alg:chaining-primary-secondary}
\small
\begin{algorithmic}[1]
  \Require chains \(C\), query intervals \(I_c\), scores \(s_c\), limits \(\lambda,L\)
  \Ensure ordered primary list \(P\) and parent \(\pi(c)\) for every chain
  \State \(P\gets[\,]\)
  \For{\(c\in C\) in decreasing order of \(s_c\)}
    \State \(u\gets\bigl|I_c\setminus\bigcup_{p\in P}I_p\bigr|,\quad \pi(c)\gets c\)
      \Comment{initially primary}
    \For{\(p\in P\) in acceptance order}
      \State \(o\gets|I_c\cap I_p|\)
      \State \(a\gets\min(|I_c|,|I_p|),\quad b\gets\max(|I_c|,|I_p|)\)
      \If{\(o>0\land o/a-u/b>\lambda\land u\leq L\)}
        \State \(\pi(c)\gets p\) \Comment{secondary alternative to \(p\)}
        \State \textbf{break}
      \EndIf
    \EndFor
    \If{\(\pi(c)=c\)}
      \State append \(c\) to \(P\)
    \EndIf
  \EndFor
\end{algorithmic}
\end{algorithm}

A separate filter retains all primaries and limits secondary alternatives.
With the defaults, a secondary must score at least \(0.8s_{\pi(c)}\) or be
within \(2k\) of its parent's score, where \(k\) is the seed length.
Exact duplicates of the parent's coordinates are removed, and at most five
secondaries survive across the list. BioseqDB follows these selection steps
and repeats them after base alignment refines the regions
(Section~\ref{sec:alignment-postprocessing}).

Primary \emph{chains} should not be confused with primary SAM records.
Minimap2 marks one surviving nonsecondary alignment as the primary SAM record.
The other nonsecondary parts of a split mapping are \emph{supplementary}.
Secondary records instead describe alternative mappings of the same query
region.
\clearpage
